\subsection{Morita equivalence and Green's theorem}\label{sec:Morita}

Crossed products of algebras provide a host of examples of dualities
which come in the form of various levels of strong and weak
equivalences of algebras, see, e.g.,~\cite{Brodzki2006}. The most
primitive form of such dualities is provided by (strong) Morita
equivalence~\cite{Rieffel1974}. A bimodule for a pair of algebras $\alg$ and $\balg$ is a
vector space $\Mcal$ which is simultaneously
 a left $\alg$-module and a right $\balg$-module, where the left action
 of $\alg$ commutes with the right action of $\balg$: $(a\cdot\xi)\cdot
 b=a\cdot (\xi\cdot b)$ for all $a\in \alg$, $b\in
 \balg$ and $\xi\in \Mcal$. If $\alg$ and $\balg$ are $C^*$-algebras,
 we say that a bimodule $\Mcal$ is an $\alg$--$\balg$ \emph{Morita
 equivalence bimodule} (or \emph{imprimitivity bimodule}) if it
is equipped with an $\alg$-valued inner product
${}_\alg\langle\,\cdot\,,\,\cdot\,\rangle$ and a $\balg$-valued inner
product $\langle\,\cdot\,,\,\cdot\,\rangle_\balg$ satisfying the
associativity condition
\[
{}_\alg\langle \psi,\phi\rangle \cdot \xi =
\psi\cdot\langle\phi,\xi\rangle_\balg ,
\]
for all $\psi,\phi,\xi\in\Mcal$, under which $\Mcal$ is complete in
the norm closures, and such that the ideal ${}_\alg\langle
\Mcal,\Mcal\rangle$ is dense in $\alg$ and $\langle
\Mcal,\Mcal\rangle_\balg$ is dense in $\balg$. The bimodule $\Mcal$ establishes a
\emph{Morita equivalence} between the algebras $\alg$ and $\balg$, and
in this case we write $\alg\ \mbf\sim_{\text{\tiny M}}\ \balg$.

Morita equivalent $C^*$-algebras have equivalent categories of
nondegenerate $*$-representations: If $\Pim_\balg\colon \balg\to{\rm B}(\hil_\balg)$ is a
representation of $\balg$ on a Hilbert space $\hil_\balg$, then we can
construct another Hilbert space
\[
\hil_\alg:=\Mcal\otimes_\balg\hil_\balg
\]
which is the quotient of the tensor product $\Mcal\otimes\hil_\balg$ by the relation $(\xi\cdot
b)\otimes\psi-\xi\otimes\Pim_\balg(b)\psi=0$ identifying the
$\balg$-actions for
$\xi\in\Mcal$, $b\in\balg$ and $\psi\in\hil_\balg$. The inner product
on $\hil_\alg$ is given by
\[
\big\langle\xi\otimes_\balg\psi\big|
\xi'\otimes_\balg\psi'\big\rangle_{\hil_\alg} :=
\big\langle\psi\big|\Pim_\balg\big(\langle\xi,\xi'\rangle_\balg
\big)\psi'\big\rangle_{\hil_\balg} ,
\]
and a representation $\Pim_\alg\colon \alg\to\rmB(\hil_\alg)$ of the algebra
$\alg$ is defined by
\[
\Pim_\alg(a)(\xi\otimes_\balg\psi) = (a\cdot\xi) \otimes_\balg\psi
\]
for $a\in\alg$ and $\xi\otimes_\balg\psi\in\hil_\alg$; this
representation is unitary equivalent to the representation
$\Pim_\balg$. Conversely, starting with a representation of $\alg$, we
can use a conjugate $\balg$--$\alg$ equivalence bimodule
$\overline{\Mcal}$ to construct a unitary equivalent representation of
$\balg$; then there are surjective bimodule homomorphisms
$\Mcal\otimes_\balg\overline{\Mcal}\to\alg$ and
$\overline{\Mcal}\otimes_\alg\Mcal\to\balg$ which satisfy a certain
transitivity law. As a particular consequence of this equivalence, Morita
equivalent algebras have homeomorphic spectra and isomorphic K-theory groups.

\begin{Example}[noncommutative two-tori]\label{ex:NCT2Morita}
A famous example of Morita equivalence in both mathematics and string
theory is provided by the
noncommutative tori ${\sf A}_\theta=\torus_\theta^2$ from
Example~\ref{ex:NCtori}. Firstly, notice from \eqref{eq:C*UV} that
changing the coset representative $\theta\in\real/\Z$ yields an identical
algebra: $\A_{\theta+m}=\A_{\theta}$ for all $m\in\Z$. Secondly, there
is an obvious
$C^*$-algebra isomorphism $\A_{-\theta}\simeq\A_\theta$ obtained by
interchanging the two generators $U$ and $V$. The converse
is also true~\cite{Rieffel1981,Rieffel1982}:
$\A_{\theta'}\simeq\A_{\theta}$ if and only if $\theta'=\theta \ {\rm
 mod}~1$. More generally, two irrational rotation $C^*$-algebras
$\A_\theta$ and $\A_{\theta'}$ are
Morita equivalent if and only if $\theta$ and $\theta'$ lie in the
same orbit under the action of $\sfGL(2,\Z)$ by fractional linear
transformations
\[
\theta'={\tt M}[\theta]:=\frac{a \theta+b}{c \theta+d} \qquad
\mbox{for} \quad {\tt M}=\bigg(\begin{matrix} a & b \\ c &
 d \end{matrix}\bigg) \in \sfGL(2,\Z) .
\]
The explicit Morita equivalence bimodules can be found
in~\cite{Rieffel1981}. On the other hand, the rational rotation
algebras $\A_\theta$ are all Morita equivalent to the commutative
algebra $C\big(\torus^2\big)$ of functions on the
two-torus~\cite{Rieffel1982}.
\end{Example}

\begin{Example}[noncommutative $d$-tori]
\label{ex:NCTdMorita}
The Morita equivalences of Example~\ref{ex:NCT2Morita} generalize to
the higher-dimensional noncommutative tori $\A_\Theta=\torus^d_\Theta$
from Example~\ref{ex:NCtorid} in the following
way~\cite{Rieffel1998}. First of all, the algebra $\A_\Theta$ is
unchanged if the matrix $\Theta$ is written in another basis of~$\Z^d$: if $B\in\sfGL(d,\Z)$ with transpose $B^{\rm t}$, then there is
a $C^*$-algebra isomorphism
$\A_{B^{\rm t}\,\Theta\,B}\simeq\A_\Theta$. More generally, consider the set of real
skew-symmetric $d{\times}d$ matrices $\Theta$ whose orbits
$\mbf M[\Theta]$ are defined for all $\mbf M\in{\sf SO}(d,d;\Z)$, where
\[
\mbf M[\Theta] = (A\,\Theta+B)\,(C\,\Theta+D)^{-1} \qquad \mbox{for} \quad
\mbf M=\left(\begin{matrix} A & B \\ C & D \end{matrix}\right) \in
{\sf SO}(d,d;\Z) ,
\]
and $A$, $B$, $C$ and $D$ are $d{\times}d$ block matrices satisfying
\[
A^{\rm t} C+C^{\rm t} A = 0 = B^{\rm t} D+D^{\rm t} B \qquad
\mbox{and} \qquad A^{\rm t} D+C^{\rm t} B=\unit_d .
\]
The set of all such matrices is dense in the space of all
skew-symmetric real $d{\times}d$ matrices, and there is a Morita
equivalence
\[
\A_{\mbf M[\Theta]} \ \mbf\sim_{\text{\tiny M}} \ \A_\Theta .
\]
However, for $d>2$ the converse is not generally true: In fact, there
are algebras $\A_\Theta$ and $\A_{\Theta'}$ that are isomorphic (and
so Morita equivalent) but for which the matrices $\Theta$ and
$\Theta'$ do not belong to the same ${\sf SO}(d,d;\Z)$
orbit~\cite{Rieffel1998}.
\end{Example}

We will also need an equivariant version of Morita equivalence in
order to show that Morita equivalent algebras induce Morita
equivalent crossed products according
to~\cite[Section~5.4]{Echterhoff2017}
\begin{Theorem}\label{equivMorita}
Let $(\alg,\sfG, \alpha)$ and $(\balg,\sfG, \beta)$ be
$C^*$-dynamical systems such that
$\alg$ and $\balg$ are Morita equivalent. Then the crossed product $C^*$-algebras
$\alg\rtimes_\alpha \sfG$ and $\balg\rtimes_\beta \sfG$ are Morita equivalent if
there exists a $\sfG$-equivariant $\alg$--$\balg$ Morita equivalence bimodule $\Mcal$,
i.e., if there is a strongly continuous action $U\colon \sfG\to \Aut(\Mcal)$ of $\sfG$ on an
$\alg$--$\balg$ Morita equivalence bimodule $\Mcal$ such that
\[
U_\gamma(a\cdot \xi)=\alpha_\gamma(a)\cdot
U_\gamma(\xi) \qquad \mbox{and} \qquad U_\gamma(\xi\cdot b)=U_\gamma(\xi)\cdot
\beta_\gamma(b) ,
\]
and
\[
{}_{\alg}\langle U_\gamma(\xi),U_\gamma(\xi')\rangle=\alpha_\gamma\big({}_\alg\langle
\xi,\xi'\rangle\big) \qquad \mbox{and} \qquad \langle
U_\gamma(\xi),U_\gamma(\xi')\rangle_\balg =\beta_\gamma\big(\langle
\xi,\xi'\rangle_\balg\big) ,
\]
for
all $\gamma\in \sfG$, $\xi,\xi'\in \Mcal$, $a\in \alg$ and $b\in \balg$.
\end{Theorem}

In this paper, our main application of Morita equivalence will involve
Green's symmetric imprimitivity theorem. Let $X$ be a locally compact space, and let
$\sfH $ and $\sfK $ be locally compact groups with commuting free and
proper actions on the right and on the left on $X$, respectively. We can lift these
actions to left actions on $C_0(X)$ by defining
$({}^h\mathtt{f})(x)=\mathtt{f}\big(h^{-1}\cdot x\big)$ and $\big({}^k\mathtt{f}\big)(x)=\mathtt{f}(x\cdot k)$
for all $\mathtt{f}\in C_0(X)$, $x\in X$, $h\in \sfH$ and $k\in \sfK $. Commutativity
of the actions of~$\sfH $ and~$\sfK $ implies that there are well-defined
induced actions of~$\sfH $ and~$\sfK $
respectively on the quotient spaces $\sfK {\setminus} X$ and $X/\sfH $, and hence
respectively on the algebras $C_0(\sfK {\setminus} X)$ and
$C_0(X/\sfH )$ which we denote~$\rt$ and~$\lt$. Green's symmetric imprimitivity theorem then reads as~\cite[Corollary~4.10]{Williams2007}
\begin{Theorem} \label{thm:Green}
There is a Morita
equivalence of $C^*$-algebras
\begin{gather}\label{eq:Greenimprim}
C_0(\sfK {\setminus} X)\rtimes_\rt \sfH \ \mbf\sim_{\text{\rm \tiny M}} \ C_0(X/\sfH )\rtimes_\lt \sfK
\end{gather}
implemented by the Morita equivalence $($or imprimitivity$)$
bimodule $\Mcal$ which is the completion of $C_{\rm c}(X)$ with the
actions
\begin{gather*}
(a\cdot\xi)(x) =\int_\sfK\,a(k,x\cdot\sfH) \xi\big(k^{-1}\cdot x\big)
\Delta_\sfK(k)^{1/2} \, \dd\mu_\sfK(k) , \\
(\xi\cdot b)(x) =\int_\sfH \xi\big(x\cdot h^{-1}\big) b\big(h,\sfK\cdot x\cdot h^{-1}\big) \Delta_\sfH(h)^{-1/2} \, \dd\mu_\sfH(h) ,
\end{gather*}
for all $x\in X$, $a\in C_{\rm c}(\sfK\times X/\sfH)$, $b\in C_{\rm
 c}(\sfH\times \sfK {\setminus} X)$ and $\xi\in C_{\rm
 c}(X)$, and the inner products
\begin{gather*}
{}_\alg\langle\xi,\xi'\rangle(k,x\cdot\sfH) = \Delta_\sfK(k)^{-1/2}
\int_\sfH \xi(x\cdot h) \overline{\xi'\big(k^{-1}\cdot x\cdot h\big)} \,
\dd\mu_\sfH(h) , \nonumber\\
\langle\xi,\xi'\rangle_\balg(h,\sfK\cdot x) =
\Delta_\sfH(h)^{-1/2} \int_\sfK \overline{\xi\big(k^{-1}\cdot x\big)}
\xi'\big(k^{-1}\cdot x\cdot h\big) \, \dd\mu_\sfK(k)  ,
\end{gather*}
for all $x\in X$, $h\in\sfH$, $k\in\sfK$ and $\xi,\xi'\in C_{\rm c}(X)$.
\end{Theorem}

Theorem~\ref{thm:Green} has several useful applications and
corollaries, see, e.g.,~\cite{Rieffel1976}. A particularly relevant
special case that we shall use below is when $\sfK $ is the trivial group,
in which case~\eqref{eq:Greenimprim} reduces to the Morita equivalence
\begin{gather*}%\label{eq:GreenKe}
C_0(X)\rtimes_\rt \sfH \ \mbf\sim_{\text{\tiny M}} \ C_0(X/\sfH ) ,
\end{gather*}
illustrating the use of crossed products in describing
quotients. This equivalence can in fact be strengthened to a stable
isomorphism~\cite{Rieffel1976}
\begin{gather}\label{eq:XsfHiso}
C_0(X)\rtimes_\rt \sfH \simeq C_0(X/\sfH ) \otimes \CK\big(\Leb^2(\sfH)\big) ,
\end{gather}
where $\CK$ denotes the algebra of
compact operators.
\begin{Example}[tori]\label{ex:Moritatori}
A particularly relevant example for us is the case $X=\real^d$ with
$\sfH=\zed^d$ acting by translations $(n,x)\mapsto x+n$ for
$n\in\zed^d$ and $x\in\real^d$, which realizes the
$d$-dimensional torus $\torus^d=\real^d/\zed^d$ as a crossed product:
\[
C\big(\torus^d\big) \ \mbf\sim_{\text{\tiny M}} \ C_0\big(\real^d\big)\rtimes_\rt \zed^d .
\]
Let us illustrate the construction explicitly. The convolution algebra
$C_{\rm c}\big(\Z^d\times\real^d\big)\!\subset\! C_0\big(\R^d\big)\rtimes_\rt\Z^d$ can be
identified with the space of sequences $f=\{f_n\}_{n\in\Z^d}$ of
functions $f_n\colon \real^d\to\complex$ with the convolution product
\[
(f\star g)_n(x) = \sum_{m\in\Z^d} f_m(x) g_{n-m}(x-m) .
\]

Consider the algebra
\[
\alg:=\big\{f\in C_0\big(\real^d,\CK\big(\ell^2\big(\zed^d\big)\big)\big) \, \big| \, f(x+m)
= U_m f(x) U_{m}^{-1}\big\} ,
\]
where $U_m$ is the unitary shift operator on $\ell^2\big(\zed^d\big)$ defined by
$(U_ma)_n=a_{n-m}$ for each $m\in\Z^d$ and $a=\{a_n\}_{n\in\Z^d}$.
Define a map ${\mit\Phi}\colon C_{\rm c}\big(\Z^d\times\real^d\big)\to
C_{\rm c}\big(\real^d,\CK\big(\ell^2\big(\Z^d\big)\big)\big)$ by
\[
\big({\mit\Phi}(f)(x)\big)_{mn} = f_{m+n}(x+n) .
\]
It is easy to see that ${\mit\Phi}(f)(x+m) =
U_m {\mit\Phi}(f)(x) U_m^{-1}$ for all $f\in C_{\rm
 c}\big(\zed^d\times\real^d\big)$, and if $a=(a_{mn})\in\alg$ then defining
$f_n:=a_{0n}$ gives ${\mit\Phi}(f)=a$. It is also easy to check that
${\mit\Phi}(f\star g) = {\mit\Phi}(f) {\mit\Phi}(g)$, and consequently
$\mit\Phi$ gives an algebra isomorphism
\[
{\mit\Phi}\colon \ C_0\big(\real^d\big)\rtimes_\rt\zed^d \xrightarrow{ \ \simeq \ } \alg .
\]

The explicit Morita equivalence bimodule is now obtained from the
completion of
\[
\Mcal = \big\{\xi\in C_{\rm c}\big(\real^d,\ell^2\big(\Z^d\big)\big) \, \big| \, \xi(x+m) =
U_m \xi(x) \big\} .
\]
The left action of the algebra $\alg={\mit\Phi}\big(C_{\rm c}\big(\Z^d\times\real^d\big)\big)$ is by left
matrix multiplication on $\Mcal$:
\begin{align*}
\alg\times\Mcal &\longrightarrow\Mcal , \\
(a,\xi) & \longmapsto a\cdot\xi , \qquad (a\cdot\xi)_n=\sum_{m\in\Z^d}
a_{nm} \xi_m ,
\end{align*}
while the right action of the algebra $C\big(\torus^d\big)$ is by right pointwise
multiplication on~$\Mcal$:
\begin{align*}
\Mcal\times C\big(\torus^d\big) & \longrightarrow\Mcal , \nonumber \\
(\xi,b) & \longmapsto \xi\cdot b , \qquad (\xi\cdot b)_n=\xi_n b .
\end{align*}
The left and right inner
products are respectively given by
\begin{gather*}
{}_\alg\langle\xi,\eta\rangle = \xi\otimes\eta^* , \\
\langle\xi,\eta\rangle_{C(\torus^d)} = \sum_{n\in\zed^d}
\overline{\xi_n} \eta_n ,
\end{gather*}
for $\xi,\eta\in\Mcal$.
Together with the isomorphism $\mit\Phi$, this establishes a
Morita equivalence between the algebras $C_0\big(\real^d\big)\rtimes_\rt\zed^d$
and $C\big(\torus^d\big)$.
\end{Example}
